\providecommand{\LeanRouteExponent}{3.58}
\providecommand{\LeanRouteExponentLo}{3.24}
\providecommand{\LeanRouteExponentHi}{3.92}
\providecommand{\LeanRouteRSquared}{0.989}
\providecommand{\LeanRoutePointCount}{9}
\providecommand{\LeanRouteBudgetMin}{10}
\providecommand{\LeanRouteThresholdKB}{28}
\providecommand{\LeanRouteThresholdBytes}{27795}
\providecommand{\LeanRouteHost}{AMD EPYC 7B13, 16 threads, 126\,GiB}
\providecommand{\LeanRouteFitA}{-30.207075}
\providecommand{\LeanRouteFitB}{3.577185}
\providecommand{\LeanRouteFitDomainLo}{919}
\providecommand{\LeanRouteFitDomainHi}{24686}
\providecommand{\LeanRouteRSqBytes}{0.989}
\providecommand{\LeanRouteRSqRows}{0.962}
\providecommand{\LeanRouteRSqNodes}{0.909}
\providecommand{\LeanRouteCensoredCount}{5}
\providecommand{\LeanRouteGlcCount}{15}
\providecommand{\LeanRouteGlcMedian}{0.36\,s}
\providecommand{\LeanRouteGlcMax}{15\,s}
\providecommand{\LeanRouteQtslabAcceptCount}{7}
\providecommand{\LeanRouteQtslabCensoredCount}{7}
\providecommand{\LeanRouteQtslabOofCount}{1}
\providecommand{\LeanRouteArmHost}{Cortex-X925 + Cortex-A725}
\providecommand{\GlcRouteExponent}{1.29}
\providecommand{\GlcRouteExponentLo}{0.75}
\providecommand{\GlcRouteExponentHi}{1.83}
\providecommand{\GlcRouteRSquared}{0.67}
\providecommand{\GlcRouteFitA}{-13.722866}
\providecommand{\GlcRouteFitB}{1.288847}
\providecommand{\GlcRouteFitDomainLo}{492}
\providecommand{\GlcRouteFitDomainHi}{126115}
\providecommand{\QtslabRouteExponent}{3.10}
\providecommand{\QtslabRouteExponentLo}{2.30}
\providecommand{\QtslabRouteExponentHi}{3.91}
\providecommand{\QtslabRouteRSquared}{0.95}
\providecommand{\QtslabRouteFitA}{-24.411397}
\providecommand{\QtslabRouteFitB}{3.101501}
\providecommand{\QtslabRouteFitDomainLo}{925}
\providecommand{\QtslabRouteFitDomainHi}{17597}


\section{The Lean route: how much Lean?}
\label{sec:exp-lean-route}

The toolkit's third checking path --- the last of the three checking
routes whose division of labor the front matter states
(\bookref{ch:exec-summary}{the Executive Summary}) --- loads the same
sealed certificate container that every other checker consumes into
Lean, reconstructs the claim, and lets Lean accept it as a theorem. The route's decode,
reconstruction, and validity elaboration are ordinary Lean definitions;
its acceptance theorem says that anything the route accepts is
semantically valid --- valid, that is, under every admissible
valuation of the hypernet the certificate declares: the same
Bernstein Logic valuation semantics the toolkit's other checkers
answer to --- and the proof is one line, because the accepted
value carries the validity proof built during elaboration. The same
definitions are consumed two ways: \emph{compiled}, as a native
executable produced by the Lean compiler, and \emph{kernel}, where a
generated receipt module asks the kernel itself to evaluate the entire
pipeline over the embedded bytes (\texttt{decide}) --- no compiler, no
extraction, just the trusted evaluator. The route exists to make a
measured statement, not a rhetorical one: the expectation has always
been that the kernel path is far too slow for machine-scale
certificates, and this section measures how far.

\subsection{The ladder}

The measurement corpus is a size ladder of fifteen production
certificates --- sealed containers from the working producer stack,
spanning $603$ bytes to $126{,}216$ bytes ($\times 209$) across seven
certificate families, each rung accepted by the independent reference
checker before any Lean timing was attempted. The ladder is committed
byte-for-byte in \path{lean-route/fixtures/}; as elsewhere in this
repository, the certificates stand on their own, and no producer is
part of the trust story. All timings in this section were taken on one
otherwise-idle machine (\LeanRouteHost), single job at a time, with a
\LeanRouteBudgetMin-minute budget per certificate per route;
Table~\ref{tab:lean-route-ladder} shows the ladder. An em-dash marks a
certificate whose run exceeded the budget; \emph{o.o.m.}\ marks one
that exhausted the machine's memory first; \emph{o.o.f.}\ (out of
fragment) marks the scope boundary discussed below --- the route
declines any certificate outside its accepted fragment rather than
guess at it, a refusal that costs completeness and never soundness ---
and kernel peaks for non-completing runs are reported as floors (the
resident size when the run was cut).

\begin{table}[htbp]
\centering\footnotesize
{\setlength{\tabcolsep}{4pt}%
\begin{tabular}{@{}lrrrrr@{}}
\toprule
certificate & bytes & proof nodes & compiled & kernel & kernel peak \\
\midrule
\texttt{enclosure\_v1\_golden\_cos6} & 603 & 1 & \emph{o.o.f.} & \emph{o.o.f.} &  \\
\texttt{ctmc\_sis\_gen} & 919 & 3 & 2.4\,ms & \emph{o.o.m.} & $\geq$124\,GiB \\
\texttt{ctmc\_fk\_mid} & 1{,}793 & 35 & 51\,ms & \emph{o.o.m.} & $\geq$124\,GiB \\
\texttt{micro2\_stable} & 2{,}184 & 96 & 0.12\,s & --- & $\geq$0.75\,GiB \\
\texttt{cartpole\_core\_decrease\_s0001} & 4{,}068 & 65 & 0.43\,s & --- & $\geq$0.76\,GiB \\
\texttt{rode\_uniform\_decrease} & 7{,}652 & 143 & 2.9\,s & --- & $\geq$0.76\,GiB \\
\texttt{cartpole\_dp\_core\_decrease} & 9{,}336 & 197 & 9.6\,s & n.a. &  \\
\texttt{dp\_composite\_syzygy} & 17{,}591 & 377 & 99\,s & n.a. &  \\
\texttt{dp\_composite\_eigenrate} & 22{,}056 & 647 & 7.4\,min & n.a. &  \\
\texttt{lti\_sampled\_flow} & 24{,}686 & 1{,}693 & 8\,min & n.a. &  \\
\texttt{acasx4\_init\_s0000} & 43{,}401 & 527 & --- & n.a. &  \\
\texttt{dp\_composite\_smallgain} & 46{,}015 & 1{,}081 & --- & n.a. &  \\
\texttt{acasx4\_promise\_closedloop} & 57{,}779 & 1{,}055 & --- & n.a. &  \\
\texttt{rode\_nest\_shell} & 119{,}217 & 4{,}629 & --- & n.a. &  \\
\texttt{rode\_hold\_flow} & 126{,}216 & 5{,}773 & --- & n.a. &  \\
\bottomrule
\end{tabular}}

\caption{The ladder. ``compiled'' is the median end-to-end wall clock
of repeated passes of the compiled route; ``kernel'' is the wall clock
of one kernel (\texttt{decide}) evaluation of the same bytes. Em-dashes
are runs cut at the \LeanRouteBudgetMin-minute budget; \emph{n.a.}\
marks kernel runs not attempted because every in-fragment kernel
attempt at smaller sizes had already ended cut or out of memory.}
\label{tab:lean-route-ladder}
\end{table}

\subsection{The scaling law}

Across the \LeanRoutePointCount{} rungs the compiled route completes
within budget, end-to-end time follows a power law in container bytes:
$t \propto B^{\LeanRouteExponent}$ (95\% CI $[\LeanRouteExponentLo,
\LeanRouteExponentHi]$, $r^2 = \LeanRouteRSquared$ on log--log least
squares). Container bytes is the best of the three natural size units as a
regressor ($r^2 = \LeanRouteRSqBytes$, against \LeanRouteRSqRows{} for
certificate rows and \LeanRouteRSqNodes{} for proof nodes) --- expected,
since the dominant cost is the reader: list-based CBOR decoding and
seam assembly in a purely functional style, chosen so that the same
definitions remain evaluable by the kernel. Nothing about the
certificate format requires a superlinear reader, as the linear-time
imported checkers on the same bytes demonstrate --- and as the
fast-path route's own fitted exponent on this same ladder now
measures (\GlcRouteExponent, its confidence interval containing
linear). Figure~\ref{fig:lean-route-scaling}
shows the measured points, the fitted law, and the budget line ---
and, since the trichotomy's claim is comparative, the same ladder
through the toolkit's other two checking routes, measured per
certificate on the Arm host
(section~\ref{sec:exp-lean-route-comparison} carries the
hardware and scope contract).

\begin{figure}[htbp]
\centering
\begin{tikzpicture}
\colorlet{routeEngine}{blue!70!black}
\colorlet{routeCertified}{green!45!black}
\colorlet{routeLean}{red!80!black}
\begin{loglogaxis}[
    width=0.82\linewidth, height=0.62\linewidth,
    xlabel={sealed container bytes},
    ylabel={end-to-end wall clock (s)},
    ymin=5e-4, ymax=5e3,
    grid=major, grid style={black!12},
    legend cell align=left, legend columns=2,
    legend style={draw=black!30, font=\small,
      at={(0.5,-0.20)}, anchor=north, /tikz/every even column/.append style={column sep=1em}},
    tick label style={font=\small},
]
\addplot[routeLean!55, thick, domain=\LeanRouteFitDomainLo:130000, samples=2]
    {exp(\LeanRouteFitA)*x^\LeanRouteFitB};
\addlegendentry{Lean fit $t \propto B^{\LeanRouteExponent}$}
\addplot[routeEngine!55, thick, domain=\GlcRouteFitDomainLo:\GlcRouteFitDomainHi, samples=2]
    {exp(\GlcRouteFitA)*x^\GlcRouteFitB};
\addlegendentry{\texttt{geolog-alpha} fit $t \propto B^{\GlcRouteExponent}$}
\addplot[routeCertified!55, thick, domain=\QtslabRouteFitDomainLo:\QtslabRouteFitDomainHi, samples=2]
    {exp(\QtslabRouteFitA)*x^\QtslabRouteFitB};
\addlegendentry{\texttt{qtslab} fit $t \propto B^{\QtslabRouteExponent}$}
\addplot[dashed, black!60, domain=550:140000, samples=2] {600};
\addlegendentry{\LeanRouteBudgetMin-minute budget}
\addplot[only marks, mark=square*, mark size=1.8pt, routeEngine]
    table {generated/lean-route-glc-points.dat};
\addlegendentry{\texttt{geolog-alpha} (fast path)}
\addplot[only marks, mark=diamond*, mark size=2.4pt, routeCertified]
    table {generated/lean-route-qtslab-points.dat};
\addlegendentry{\texttt{qtslab} (certified)}
\ifnum\LeanRouteQtslabCensoredCount>0
\addplot[only marks, mark=diamond, mark size=2.8pt, routeCertified]
    table {generated/lean-route-qtslab-censored.dat};
\addlegendentry{\texttt{qtslab}, cut at budget}
\fi
\addplot[only marks, mark=*, mark size=1.7pt, routeLean]
    table {generated/lean-route-points.dat};
\addlegendentry{Lean (compiled reflection)}
\ifnum\LeanRouteCensoredCount>0
\addplot[only marks, mark=o, mark size=2.2pt, routeLean]
    table {generated/lean-route-censored.dat};
\addlegendentry{Lean, cut at budget}
\fi
\end{loglogaxis}
\end{tikzpicture}
\caption{One ladder, three checking routes: the fast-path engine
(\texttt{geolog-alpha}, blue squares, on the staged-chain sealed
containers), the certified checker (\texttt{qtslab}, green diamonds,
on the same certificates re-marked content-resolved --- a
container-table transform under which every certificate row is
byte-identical), and the compiled Lean reflection route (red circles)
with its fitted law. Blue and green medians are one process, end to
end, pinned to one core of the Arm host; the red points are this
section's own protocol on its dedicated machine (the footnote of
section~\ref{sec:exp-lean-route-comparison} carries both). Open marks
of either color sit \emph{at} the dashed
budget line: each true time lies above its mark. A rung with no mark
on some route is a scope datum, not a failure --- that route classes
the certificate outside its accepted fragment and declines it (the
caption of Table~\ref{tab:lean-route-ladder} and
section~\ref{sec:exp-lean-route-comparison}'s footnote carry the
per-route scope and hardware contract). The Lean fit uses only
uncensored red points; where it crosses the budget ---
$\approx\LeanRouteThresholdKB$\,KB of sealed container --- is the
admission threshold of Table~\ref{tab:lean-route-admission}. Each
solid line is the same log--log fit to its route's accepted medians:
\texttt{geolog-alpha} at exponent \GlcRouteExponent{} (95\% CI
$[\GlcRouteExponentLo,\GlcRouteExponentHi]$, containing linear;
$r^2=\GlcRouteRSquared$), \texttt{qtslab} at \QtslabRouteExponent{}
($[\QtslabRouteExponentLo,\QtslabRouteExponentHi]$,
$r^2=\QtslabRouteRSquared$), Lean at \LeanRouteExponent{}
($[\LeanRouteExponentLo,\LeanRouteExponentHi]$) --- one near-linear
route and two superlinear ones, measured on the same certificates.}
\label{fig:lean-route-scaling}
\end{figure}

The kernel mode needs no fitted curve to make its point. Its smallest
full acceptance during the route's development --- the $919$-byte
certificate, second rung of the ladder --- cost $965$ seconds and a
$233$\,GiB peak on a $251$\,GiB development machine (the toolkit
repository's route documentation carries that receipt); the memory is
the kernel's reduction state, not the certificate. On this section's
$126$\,GiB reference machine the same certificate exhausts neither
budget nor patience but memory, and every in-fragment kernel attempt
on the ladder ended cut at the budget, out of memory, or predictably
unattempted (Table~\ref{tab:lean-route-ladder}). The failure mode is
content-dependent, not size-monotone: the $919$-byte certificate took
the whole machine's memory in under eight minutes, while certificates
two to eight times its size ground past the budget inside a
gigabyte --- a $\times 160$ spread in peak residency between
neighboring rungs. What a given certificate makes the kernel's
reduction state do is a property of its content (the working
hypothesis is exact-rational arithmetic in kernel terms, untested
here), not of its byte count. A certificate of under a kilobyte
that the compiled route accepts in milliseconds is already past what
the kernel evaluator will do inside a ten-minute, hundred-gigabyte
envelope; that is the measured content of ``too slow at machine
scale,'' stated with its numbers attached.

\subsection{The fragment boundary}

One ladder rung is not a timing datum but a scope datum: the smallest
certificate ($603$ bytes, an interval-enclosure certificate whose
obligations are discharged by \emph{primitive applications}) is out of
fragment for the route --- the table's \emph{o.o.f.} --- and the route
declines it: compiled and kernel agree, and the kernel's refusal is
itself a theorem (\texttt{decide} proves the acceptance predicate
false on those bytes). The same bytes are accepted by the reference
checker, which settles the direction of the disagreement: the route's
accepted fragment is the first-order pinned-vocabulary rule set,
primitive applications lie outside it by design, and a checker asked
about a certificate it cannot speak for declines it rather than
guess. Every checker in this chapter draws some such boundary, always
with the same sign --- an out-of-fragment refusal costs completeness,
never soundness; the certified checker classes this same rung outside
\emph{its} fragment too, naming the reason in its verdict
(\texttt{fun\_outside\_fragment}), while the fast-path lane discharges
it (Figure~\ref{fig:lean-route-scaling}). Admission to the Lean route
is therefore a conjunction: small enough, \emph{and} inside the
fragment. A species' first Lean attempt should be its smallest
instance, which settles fragment membership in milliseconds before
any batch of timings is scheduled.

\subsection{The admission gate}

Table~\ref{tab:lean-route-admission} applies the fitted law to the
certificate species of this chapter's campaigns: per species, the
byte range of its committed certificates (container overhead added),
the predicted compiled-route time at the median and maximum, and the
verdict against the \LeanRouteBudgetMin-minute budget. The table is the
gate the re-run campaigns consume: a Lean-reflection timing column is
attempted only for species that pass (or, at the runner's discretion,
the passing fraction of a \emph{partial} species), and the kernel mode
is attempted for none --- the measured points above place its envelope
below every production species' floor.

\begin{table}[htbp]
\centering\footnotesize
{\setlength{\tabcolsep}{4pt}%
\begin{tabular}{@{}p{10.5em}rllll@{}}
\toprule
species & max B & fragment & pred.\ med. & pred.\ max & verdict \\
\midrule
hardware liveness (circuits) & 2{,}627 & untested & 9.1\,ms & 0.22\,s & pass \\
Atari certified replay & 5{,}113 & untested & 22\,ms & 1.9\,s & pass \\
FrozenLake 4x4 (det + slippery + sr34) & 52{,}015 & exp.\ in & 15\,s & 97\,min & partial \\
CliffWalking & 71{,}872 & exp.\ in & 14\,min & 5.1\,h & fail \\
FrozenLake 8x8 & 500{,}289 & exp.\ in & 5.5\,h & 0.59\,yr & fail \\
Blackjack & 115{,}779 & exp.\ in & 10\,h & 28\,h & fail \\
Taxi (all weather/passenger variants) & 3{,}102{,}602 & exp.\ in & 7\,yr & 400\,yr & fail \\
LunarLander / LunarLanderContinuous (statistical tier) & 12{,}226{,}428 & untested & $\sim$$10^{4}$\,yr & $\sim$$10^{5}$\,yr & fail \\
published-table supermartingale corpus & 50{,}000 & exp.\ in & 97\,ms & 84\,min & partial \\
semiconductor fluid ladder (minifab) & 18{,}451 & untested & 65\,s & 2.5\,min & pass \\
semiconductor fluid ladder (mimac) & 401{,}455 & untested & 0.18\,yr & 0.27\,yr & fail \\
interval-enclosure certificates (primitive applications) & 603 & out & 4.4\,ms & 4.4\,ms & \emph{out of fragment} \\
\bottomrule
\end{tabular}}

\caption{The admission gate. Predictions apply the fitted law to each
species' certificate sizes in sealed-container bytes (census rows
stated in committed-artifact bytes carry the measured container
overhead as a declared correction); beyond the fitted range
($\approx 25$\,KB) they are
extrapolations --- one-sidedly corroborated by the budget-censored
rungs up to $126$\,KB, order-of-magnitude statements past that.
``Fragment'' marks measured membership only for the ladder's own
species; \emph{untested} species settle theirs with a one-instance
probe before any batch attempt.}
\label{tab:lean-route-admission}
\end{table}

\subsection{The measurement before this one}
\label{sec:exp-lean-route-history}

This ladder is the route's second measurement, not its first. In July
2026, on the toolkit's predecessor binary certificate format (since
retired), the same question was measured by the other road into the
kernel --- \emph{elaboration}: an emitter translated certificate bytes
into generated Lean source ($\times 8$ expansion, byte-stable per-node
templates across certificate kinds), and the elaborator produced
kernel-checked object files (a further $\times 24$--$25$) closing with
the same species of validity theorem. Three production certificates
were taken to full kernel verdicts that way --- $197$, $377$, and
$1{,}081$ proof nodes, at $3.4$, $9$--$13$, and $28$ minutes of
elaboration respectively, against $3$--$10$ \emph{seconds} for the
certified checker build of the same era on identical certificate
content and comparable hardware: a measured factor of $\times 30$ to
$\times 250$ at kilobyte scale, growing with size. The same campaign
took a streaming census of a real machine-scale certificate --- a
$43{,}285{,}639$-node, $1.24$\,GB closed-loop neural-network
verification certificate, semantically re-checked by a reference
checker in $2.4$ hours --- and projected the elaboration route on it:
$\approx 7$\,GB of generated Lean source, $\approx 66$ million
declarations in one namespace (about one hundred mathlibs),
$\approx 176$\,GB of object files, and $1.4$--$2.7$ core-years of
serial elaboration, with memory the first wall. The numbers, the
census, and the projection's stated assumptions are recorded in
\path{lean-route/HISTORY.md}. That projection is why the present route
abandons source generation altogether: a fixed decoder and checker are
defined once, certificates arrive as ordinary input bytes, and the
measured question becomes the reader's scaling law above rather than
gigabyte-scale generated source.

\subsection{The route beside the other two}
\label{sec:exp-lean-route-comparison}

The trichotomy's claim is comparative, so the honest summary is
too.\footnote{Hardware, per the conventions: the certified checker is
timed on one core of an Arm host --- the ISA its soundness theorem
grounds in --- with its own cross-machine spread measured under
$\times 1.15$; in Figure~\ref{fig:lean-route-scaling} the two
production routes' ladder points are one process per certificate,
pinned to one performance core of one Arm box (\LeanRouteArmHost),
while the red points remain this section's reflection timings
(\LeanRouteHost); the elaboration-era numbers are as recorded in
\path{lean-route/HISTORY.md}. Every contrast quoted here is two or
more orders of magnitude, and the measured cross-machine factor on the
compiled route is a factor of a few --- hardware choice cannot close
any of them. Container forms, per route: the fast path consumes the
staged-chain re-seals (\path{lean-route/GLC-CERT.md}), the certified
checker the content-resolved mints
(\path{lean-route/tools/mint-qtslab-vintage.py} states the
container-table-only transform), and both carry the same certificate
rows byte-for-byte as the committed fixtures the Lean route reads.
Scope, shared by all three routes on this corpus: these
certificates natively encode the Bernstein Logic fragment, which is
the only fragment the certified checker consumes
(\path{theories/BernsteinLogic.geolog}); the full QTSL theory is
checked on the engine-side container lanes of
section~\ref{sec:exp-repro}.} On this chapter's stored supermartingale corpus ---
$4{,}105$ sealed containers, each one verification condition's
Bernstein Logic discharge at roughly one to fifty kilobytes --- the
certified checker's current build replays every container at a
pooled median of $30$\,ms, maximum $117$\,ms, on one Arm core
(Table~\ref{tab:exp-certified-checker}), and the fast-path engine's
sealed-container lane re-checks the whole corpus in wall-clock seconds
(\path{supermartingales/GLC-CERT.md}). The compiled reflection route
has not been batch-run on that corpus; the admission gate prices it
(Table~\ref{tab:lean-route-admission}): verdict \emph{partial}, with a
predicted median near $0.1$\,s and a predicted maximum near $84$
minutes, so only the passing fraction would ever be attempted.
Figure~\ref{fig:lean-route-scaling} makes the comparison direct on
the ladder itself, every route in its current build on identical
certificate content: the fast path near-linear in bytes (exponent
\GlcRouteExponent), the reflection reader at $B^{\LeanRouteExponent}$,
and the certified build accepting \LeanRouteQtslabAcceptCount{} rungs
and cut at the budget on \LeanRouteQtslabCensoredCount. Within the
ladder the certified wall grows as $B^{\QtslabRouteExponent}$
($r^2=\QtslabRouteRSquared$) --- statistically indistinguishable from
the reflection reader's exponent, while the fast path's confidence
interval excludes both --- and across corpora content dominates
bytes: the stored corpus above tops out at $117$\,ms per certificate
at sizes where the ladder's composite certificates cost minutes,
which is how a $30$\,ms pooled median there and these walls here are
both true. A
species aiming certificates at the certified route should therefore
run the same one-instance probe the admission gate prescribes for
Lean: fragment membership and cost are both content-properties,
settled cheaply on the smallest instance. The historical cross-check
stands beside the figure: on the composite small-gain rung
($1{,}081$ proof nodes, $46$\,KB sealed), the certified checker of
the elaboration era took about ten seconds on the same content in its
own format's bytes (today's build, reading the unified sealed
container, is cut at the budget on that rung), against past the
ten-minute budget through the compiled reflection route and
twenty-eight minutes through the retired elaboration route.

The conclusion the measurements support, stated as the front matter's
expectation deserves: \emph{confirmed, with its boundary measured}.
For the kernel evaluator the confirmation is categorical --- a
certificate under a kilobyte already costs sixteen minutes and a
quarter-terabyte of reduction state, and nothing larger completes at
all --- and for the compiled route it is a statement about scaling
rather than about the route: below a few kilobytes the compiled
reflection path is genuinely fast (the $919$-byte rung's $2.4$\,ms is
\emph{under} the certified build's corpus median, on different
certificates of the same format), but at $\times B^{3.6}$ against the
production checkers' linear-with-content growth, the ratio passes two
orders of magnitude by fifty kilobytes and machine scale is out of
reach by shape, not by constant factors. The route earns its keep on
the other side of the ledger: full kernel verdicts on small
certificates tie the certificate semantics to a proof kernel
scrutinized by a community that owes this toolkit nothing --- and the
measured boundary is precisely why the certified binary, not the
kernel, walks the forty-million-node instances.

\paragraph{What is not claimed.} The law is a property of this route's
current reader on one machine class, not of Lean, whose kernel and
compiler are here being used far outside their engineered envelope on
purpose; an index-structured reader would move the compiled curve
substantially, and nothing in the trichotomy's argument depends on it
not moving. Cross-machine transfer of the constants is measured to be
worth a factor of a few (the $2{,}184$-byte rung runs $27.6$\,ms on the
development machine cited above and $\approx 120$\,ms here); the
exponent, not the coefficient, is the portable content. Species
predictions inherit a byte-basis caveat until each campaign banks
per-instance container sizes, and the admission table regenerates from
those receipts by the same make target that produced it. The
elaboration-era numbers of \path{lean-route/HISTORY.md} are
development-record measurements on a retired format, quoted for the
route's history with their measured/projected distinctions preserved;
they are not re-derivable from this repository, and no other number in
this chapter rests on them. Finally, the ladder's timings were taken
against the committed fixtures; the same certificate content
re-sealed on the unified certificate schema is already committed
beside them (\path{lean-route/out/glc-baked/}, with the
content-identity argument stated in \path{lean-route/GLC-CERT.md}),
and when the ladder re-runs on those containers this section's
numbers move with them, per the reproduction discipline of
section~\ref{sec:exp-repro}.
